% Reviewed translation: PDF pages 342--345 / printed pages 328--331.
\MathBookSourcePage{342}
Theorem~\ref{thm:incompressible-decoupled-stream-pressure} 已经证明, 当
$\boldsymbol f\in L^r(\Omega)^2$,
$\operatorname{curl}\boldsymbol u\in W^{1,r}(\Omega)$ 且
$p\in W^{1,r}(\Omega)$ 时, Stokes Problem~
\eqref{eq:navier-centered-stokes-operator} 等价于:

求 $\psi\in\Phi_s$ 和 $\omega\in W^{1,r}(\Omega)$, 使得
\begin{subequations}\label{eq:navier-centered-stream-stokes}
\begin{align}
\left\{
\begin{aligned}
(\operatorname{curl}\omega,\operatorname{curl}\phi)
&=(\boldsymbol f,\operatorname{curl}\phi)
&&\forall\phi\in\Phi_s,\\
(\operatorname{curl}\psi,\operatorname{curl}\mu)
&=(\omega,\mu)
&&\forall\mu\in W^{1,r}(\Omega),
\end{aligned}
\right.
\label{eq:navier-centered-stream-stokes-vorticity}\\
\intertext{求 $p\in W^{1,r}(\Omega)\cap L_0^2(\Omega)$, 使得}
(\operatorname{grad}p,\operatorname{grad}q)
&=(\boldsymbol f-\operatorname{curl}\omega,
\operatorname{grad}q)
\qquad\forall q\in W^{1,s}(\Omega).
\label{eq:navier-centered-stream-stokes-pressure}
\end{align}
\end{subequations}
这里 $\boldsymbol u=\operatorname{curl}\psi$ 且
$\omega=\operatorname{curl}\boldsymbol u$.

这里仍需对每个右端 $\boldsymbol f\in L^r(\Omega)^2$ 把 Stokes 问题写成
\eqref{eq:navier-centered-stream-stokes} 的形式. 因此假设
\eqref{eq:navier-centered-gp-regularity} 成立. 置
\[
X=\{\operatorname{curl}\phi;\ \phi\in\Phi_s\}
\times L^2(\Omega)\times L_0^2(\Omega),
\qquad Y=L^r(\Omega)^2,
\]
则 Stokes 算子 $T$ 由
\[
T\in\mathcal L(Y;X),
\qquad
T\boldsymbol f=(\operatorname{curl}\psi,\omega,p),
\quad\text{为 \eqref{eq:navier-centered-stream-stokes} 的解}
\]
定义.

接着, 由 \eqref{eq:navier-steady-two-dimensional-component-identities},
对流项满足恒等式
\begin{equation}\label{eq:navier-centered-stream-convection-curl}
\begin{aligned}
\int_\Omega\left(\sum_{j=1}^{2}u_j
\frac{\partial\boldsymbol u}{\partial x_j}\right)
\cdot\operatorname{curl}\phi\,dx
&=\int_\Omega\omega\operatorname{grad}\psi
\cdot\operatorname{curl}\phi\,dx\\
&=\int_\Omega\omega\left(
\frac{\partial\psi}{\partial x_1}
\frac{\partial\phi}{\partial x_2}
-\frac{\partial\psi}{\partial x_2}
\frac{\partial\phi}{\partial x_1}
\right)dx,
\end{aligned}
\end{equation}
以及
\begin{equation}\label{eq:navier-centered-stream-convection-gradient}
\begin{aligned}
\int_\Omega\left(\sum_{j=1}^{2}u_j
\frac{\partial\boldsymbol u}{\partial x_j}\right)
\cdot\operatorname{grad}q\,dx
={}&\int_\Omega\omega\operatorname{grad}\psi
\cdot\operatorname{grad}q\,dx\\
&+\frac12\int_\Omega
\operatorname{grad}(\|\boldsymbol u\|^2)
\cdot\operatorname{grad}q\,dx,
\end{aligned}
\end{equation}
其中 $\operatorname{curl}\psi=\boldsymbol u$,
$\omega=\operatorname{curl}\boldsymbol u$, $\|\cdot\|$ 表示 Euclidean 范数.
因此用映射
\begin{equation}\label{eq:navier-centered-stream-nonlinearity}
G(\lambda,u)=\lambda(\omega\operatorname{grad}\psi-\boldsymbol f),
\qquad
\lambda\in\mathbb R_+,
\quad
u=(\operatorname{curl}\psi,\omega,p)\in X
\end{equation}
引入非线性, 并约定把项 $\frac12\|\boldsymbol u\|^2$ 纳入压力,
也就是说改用运动压力
$p^*=p+\frac12\|\boldsymbol u\|^2$.
因为 $s\geq4$, 项 $\omega\operatorname{grad}\psi$ 和
$\|\boldsymbol u\|^2$ 分别属于 $L^{4/3}(\Omega)^2$ 和 $L^2(\Omega)$.

现在对 $\boldsymbol f\in L^r(\Omega)^2$ 考虑下列问题:
求 $\psi\in\Phi_s$ 和 $\omega\in W^{1,r}(\Omega)$, 使得
\begin{subequations}\label{eq:navier-centered-stream-continuous}
\begin{align}
\left\{
\begin{aligned}
\nu(\operatorname{curl}\omega,\operatorname{curl}\phi)
+(\omega\operatorname{grad}\psi,\operatorname{curl}\phi)
&=(\boldsymbol f,\operatorname{curl}\phi)
&&\forall\phi\in\Phi_s,\\
(\operatorname{curl}\psi,\operatorname{curl}\mu)
&=(\omega,\mu)
&&\forall\mu\in W^{1,r}(\Omega),
\end{aligned}
\right.
\label{eq:navier-centered-stream-continuous-vorticity}
\end{align}
求 $p^*\in W^{1,r}(\Omega)\cap L_0^2(\Omega)$, 使得
\begin{align}
(\operatorname{grad}p^*,\operatorname{grad}q)
&=(\boldsymbol f-\omega\operatorname{grad}\psi
-\nu\operatorname{curl}\omega,\operatorname{grad}q)\notag\\
&\hspace{8em}\forall q\in W^{1,s}(\Omega).
\label{eq:navier-centered-stream-continuous-pressure}
\end{align}
\end{subequations}

\MathBookSourcePage{343}
采用上面的记号, 该问题写成
\[
F(\lambda,u)\equiv u+TG(\lambda,u)=0,
\]
其中 $\lambda=1/\nu$,
$u=(\operatorname{curl}\psi,\omega,\lambda p^*)\in X$,
$G$ 由 \eqref{eq:navier-centered-stream-nonlinearity} 定义,
$T$ 由 \eqref{eq:navier-centered-stream-stokes} 定义.

再次通过直接计算可知: 若 $(\boldsymbol u,p)$ 是 Problem~
\eqref{eq:navier-centered-continuous-mixed} 的解,
$\operatorname{curl}\boldsymbol u,p\in W^{1,r}(\Omega)$ 且
$\boldsymbol f\in L^r(\Omega)^2$, 则满足
\begin{equation}\label{eq:navier-centered-stream-variable-relation}
\operatorname{curl}\psi=\boldsymbol u,
\qquad
\operatorname{curl}\boldsymbol u=\omega,
\qquad
p^*=p+\frac12\|\boldsymbol u\|^2
\end{equation}
的三元组 $(\operatorname{curl}\psi,\omega,p^*)$ 是 Problem~
\eqref{eq:navier-centered-stream-continuous} 的解. 反之,
\eqref{eq:navier-centered-stream-continuous} 的每个解
$(\operatorname{curl}\psi,\omega,p^*)$ 都满足 $\omega=-\Delta\psi$,
且由 \eqref{eq:navier-centered-stream-variable-relation} 定义的
$(\boldsymbol u,p)$ 满足 \eqref{eq:navier-centered-continuous-mixed}.
此外, 当 Stokes 算子具有正则性
\eqref{eq:navier-centered-gp-regularity} 时, Problem~
\eqref{eq:navier-centered-stream-continuous} 对某个实数
$\gamma\in[r,2]$ 和右端 $\boldsymbol f\in L^\gamma(\Omega)^2$ 的每个解
$u=(\operatorname{curl}\psi,\omega,p^*)$ 都具有正则性
\[
\psi\in W^{3,\gamma}(\Omega),
\qquad
\omega\in W^{1,\gamma}(\Omega),
\qquad
p^*\in W^{1,\gamma}(\Omega).
\]
再者, 正如 Remark~\ref{rem:navier-centered-gp-nonsingular-equivalence},
Problem~\eqref{eq:navier-centered-continuous-mixed} 对右端
$\boldsymbol f\in L^r(\Omega)^2$ 的每个非奇异解分支也是
Problem~\eqref{eq:navier-centered-stream-continuous} 的非奇异解分支,
反之亦然.

关于逼近, 为了能完全剖分 $\Omega$, 假设 $\Omega$ 是
$\mathbb R^2$ 中的多边形区域. 令 $\mathcal T_h$ 是 $\bar\Omega$ 的一族
三角剖分, $l\geq1$ 是固定整数. 取
\begin{equation}\label{eq:navier-centered-stream-discrete-spaces}
\left\{
\begin{aligned}
\Theta_h
&=\{\theta_h\in\mathcal C^0(\bar\Omega);
\ \theta_h|_\kappa\in P_l
\quad\forall\kappa\in\mathcal T_h\}
\subset W^{1,\infty}(\Omega),\\
\Phi_h
&=\Theta_h\cap\Phi
=\{\phi_h\in\Theta_h;
\ \phi_h|_{\Gamma_0}=0,
\ \phi_h|_{\Gamma_i}=\text{an arbitrary constant},
\ 1\leq i\leq p\},\\
Q_h
&=\{q_h\in\mathcal C^0(\bar\Omega)\cap L_0^2(\Omega);
\ q_h|_\kappa\in P_k
\quad\forall\kappa\in\mathcal T_h\},
\qquad k=\min(1,l-1),\\
X_h&=\{\operatorname{curl}\phi_h;\ \phi_h\in\Phi_h\}
\times\Theta_h\times Q_h\subset X.
\end{aligned}
\right.
\end{equation}
利用这些空间, Stokes 问题逼近如下:
求 $\psi_h\in\Phi_h$ 和 $\omega_h\in\Theta_h$, 使得
\begin{subequations}\label{eq:navier-centered-stream-discrete-stokes}
\begin{align}
\left\{
\begin{aligned}
(\operatorname{curl}\omega_h,\operatorname{curl}\phi_h)
&=(\boldsymbol f,\operatorname{curl}\phi_h)
&&\forall\phi_h\in\Phi_h,\\
(\operatorname{curl}\psi_h,\operatorname{curl}\mu_h)
&=(\omega_h,\mu_h)
&&\forall\mu_h\in\Theta_h,
\end{aligned}
\right.
\label{eq:navier-centered-stream-discrete-stokes-vorticity}\\
\intertext{求 $p_h\in Q_h$, 使得}
(\operatorname{grad}p_h,\operatorname{grad}q_h)
&=(\boldsymbol f-\operatorname{curl}\omega_h,
\operatorname{grad}q_h)
\qquad\forall q_h\in Q_h.
\label{eq:navier-centered-stream-discrete-stokes-pressure}
\end{align}
\end{subequations}
相应的算子 $T_h\in\mathcal L(Y;X_h)$ 由
\[
T_h\boldsymbol f=(\operatorname{curl}\psi_h,\omega_h,p_h),
\quad\text{为 \eqref{eq:navier-centered-stream-discrete-stokes} 的解}
\]
定义. 类似地, Navier--Stokes 问题
\eqref{eq:navier-centered-stream-continuous} 离散化为:

\MathBookSourcePage{344}
求 $\psi_h\in\Phi_h$ 和 $\omega_h\in\Theta_h$, 使得
\begin{subequations}\label{eq:navier-centered-stream-discrete}
\begin{align}
\left\{
\begin{aligned}
\nu(\operatorname{curl}\omega_h,\operatorname{curl}\phi_h)
+(\omega_h\operatorname{grad}\psi_h,
\operatorname{curl}\phi_h)
&=(\boldsymbol f,\operatorname{curl}\phi_h)
&&\forall\phi_h\in\Phi_h,\\
(\operatorname{curl}\psi_h,\operatorname{curl}\mu_h)
&=(\omega_h,\mu_h)
&&\forall\mu_h\in\Theta_h,
\end{aligned}
\right.
\label{eq:navier-centered-stream-discrete-vorticity}
\end{align}
求 $p_h^*\in Q_h$, 使得
\begin{align}
(\operatorname{grad}p_h^*,\operatorname{grad}q_h)
&=(\boldsymbol f-\omega_h\operatorname{grad}\psi_h
-\nu\operatorname{curl}\omega_h,
\operatorname{grad}q_h)\notag\\
&\hspace{8em}\forall q_h\in Q_h.
\label{eq:navier-centered-stream-discrete-pressure}
\end{align}
\end{subequations}
换言之, 该问题也可写成
\[
F_h(\lambda,u_h)\equiv u_h+T_hG(\lambda,u_h)=0,
\]
其中 $\lambda=1/\nu$,
$u_h=(\operatorname{curl}\psi_h,\omega_h,\lambda p_h^*)\in X_h$,
$G$ 由 \eqref{eq:navier-centered-stream-nonlinearity} 定义,
$T_h$ 由 \eqref{eq:navier-centered-stream-discrete-stokes} 定义.

现在以 $Z=Y$ 应用 Theorem~\ref{thm:navier-branch-operator-approximation}.
回忆 Subsection~\ref{subsec:incompressible-further-error-refinement} 中得到的算子
$T_h$ 的逼近性质(参见 Theorem~
\ref{thm:incompressible-further-w3alpha-error}). 当 $\Omega$ 是凸多边形时,
Stokes 问题具有正则性 \eqref{eq:navier-centered-gp-regularity}; 换言之,
若 $\boldsymbol f\in L^t(\Omega)^2$, $r\leq t\leq2$, 则 Stokes 问题的解
$(\operatorname{curl}\psi,\omega,p)$ 属于
$W^{2,t}(\Omega)^2\times W^{1,t}(\Omega)\times W^{1,t}(\Omega)$, 且
\[
\|\psi\|_{3,t,\Omega}
+\|\omega\|_{1,t,\Omega}
+|p|_{1,t,\Omega}
\leq C_1\|\boldsymbol f\|_{0,t,\Omega}.
\]
因此, 若 $\mathcal T_h$ 是 $\bar\Omega$ 的一族一致正则三角剖分,
则 Problem~\eqref{eq:navier-centered-stream-discrete-stokes} 的解
$(\operatorname{curl}\psi_h,\omega_h,p_h)$ 满足
\[
|\psi-\psi_h|_{1,s,\Omega}
+\|\omega-\omega_h\|_{0,\Omega}
+\|p-p_h\|_{0,\Omega}
\leq C_2h^\alpha\|\boldsymbol f\|_{0,t,\Omega},
\]
其中 $r\leq t<2$, $1/\gamma+1/t=1$;
当 $l=1$ 时 $\alpha=1/\gamma$, 当 $l\geq2$ 时 $\alpha=2/\gamma$.
这就建立了 \eqref{eq:navier-branch-strong-operator-convergence} 和
\eqref{eq:navier-branch-z-operator-convergence}. 因而
Theorem~\ref{thm:navier-branch-operator-approximation} 的结论成立;
与 Theorem~\ref{thm:incompressible-further-smooth-error} 合并得到下列结果.

\setcounter{statement}{3}
\begin{Theorem}\label{thm:navier-centered-stream-branch}
$1^\circ)$ 设 $\Omega$ 是有界凸多边形,
$\mathcal T_h$ 是 $\bar\Omega$ 的一族一致正则三角剖分. 对
$\boldsymbol f\in L^t(\Omega)^2$, $t\in[r,2]$, 设
\[
\{(\lambda,(\operatorname{curl}\psi(\lambda),\omega(\lambda),
\lambda p^*(\lambda)));
\ \lambda=1/\nu\in\Lambda\}
\]
是 Navier--Stokes Problem~
\eqref{eq:navier-centered-stream-continuous} 的一个非奇异解分支.
则当 $h\leq h_0$ 充分小时, 存在 Problem~
\eqref{eq:navier-centered-stream-discrete} 的唯一 $C^\infty$ 解分支
\[
\{(\lambda,(\operatorname{curl}\psi_h(\lambda),\omega_h(\lambda),
\lambda p_h^*(\lambda)));
\ \lambda=1/\nu\in\Lambda\},
\]
满足
\begin{equation}\label{eq:navier-centered-stream-branch-error}
\begin{aligned}
\sup_{\lambda\in\Lambda}
\{\,&|\psi_h(\lambda)-\psi(\lambda)|_{1,s,\Omega}
+\|\omega_h(\lambda)-\omega(\lambda)\|_{0,\Omega}\\
&+\|p_h^*(\lambda)-p^*(\lambda)\|_{0,\Omega}\}
\leq Ch^\alpha,
\qquad
\alpha=\begin{cases}
1/\gamma,&l=1,\\
2/\gamma,&l\geq2,
\end{cases}
\end{aligned}
\end{equation}
其中 $1/t+1/\gamma=1$, 常数 $C$ 与 $h$ 和 $\lambda$ 无关.
当 $t=2$ 且 $l\geq2$ 或 $\psi\in W^{2,\infty}(\Omega)$ 时,
该界仍然成立. 当 $l=1$ 且 $t=2$ 时,
\eqref{eq:navier-centered-stream-branch-error} 的左端由
\[
C(\varepsilon)h^{1/2-\varepsilon}
\qquad\forall\varepsilon>0
\]
控制.

\MathBookSourcePage{345}
$2^\circ)$ 此外, 若对某个实数 $m\in[1,l-1/2]$, 映射
$\lambda\mapsto(\psi(\lambda),p^*(\lambda))$ 从 $\Lambda$ 到
\[
[H^{m+2}(\Omega)\cap W^{m+3/2,\infty}(\Omega)]
\times H^m(\Omega)
\]
连续, 则对所有 $\lambda\in\Lambda$ 有误差估计
\begin{equation}\label{eq:navier-centered-stream-smooth-error}
|\psi_h(\lambda)-\psi(\lambda)|_{1,s,\Omega}
+\|\omega_h(\lambda)-\omega(\lambda)\|_{0,\Omega}
+\|p_h^*(\lambda)-p^*(\lambda)\|_{0,\Omega}
\leq Kh^m.
\end{equation}
\end{Theorem}

与线性情形一样, 可以在 $H_0^1$ 范数中改进上面对
$\psi_h-\psi$ 的估计. 然而, Theorem~
\ref{thm:navier-branch-sharper-norm-error} 的论证似乎不能在这里应用,
因为它必须同时作用于 $\psi$ 和 $\omega$, 而 $\omega$ 的估计不大可能改进.
我们改为引入一个更直接的对偶性论证. 因为不关心压力, 取
\[
X=\{\operatorname{curl}\phi;\ \phi\in\Phi_s\}\times L^2(\Omega),
\qquad
H=\{\operatorname{curl}\phi;\ \phi\in\Phi\}\times L^2(\Omega),
\]
并假设
\[
\lambda\longmapsto u(\lambda)
=(\operatorname{curl}\psi(\lambda),\omega(\lambda))
\]
是右端 $\boldsymbol f\in L^2(\Omega)^2$ 的 Problem~
\eqref{eq:navier-centered-stream-continuous-vorticity} 的一个非奇异解分支.
仍假设 \eqref{eq:navier-centered-gp-regularity} 成立, 从而
$u(\lambda)\in H^2(\Omega)\times H^1(\Omega)$. 为简化记号, 置
\[
D=D_uG(\lambda,u(\lambda)).
\]
回忆
\[
D\cdot v
=\lambda\{\omega(\lambda)\operatorname{grad}\phi
+\theta\operatorname{grad}\psi(\lambda)\}
\qquad
\forall v=(\operatorname{curl}\phi,\theta)\in X.
\]
引入由下式定义的算子 $D^*\in\mathcal L(H;X')$:
\begin{equation}\label{eq:navier-centered-stream-adjoint-derivative}
\langle D^*z,v\rangle
=(\operatorname{curl}\chi,D\cdot v)
\qquad
\forall z=(\operatorname{curl}\chi,\mu)\in H.
\end{equation}
为联系 $D$ 和 $D^*$, 回忆与 $X$ 相关的空间 $V$ 为
\[
V=\{v=(\operatorname{curl}\phi,\theta)\in X;
\ \widetilde b(v,\mu)=0
\quad\forall\mu\in W^{1,r}(\Omega)\},
\]
其中
\[
\widetilde b(v,\mu)
=(\operatorname{curl}\phi,\operatorname{curl}\mu)-(\theta,\mu).
\]
还回忆, 对标量积
\[
\widetilde a(u,v)=(\omega,\theta)
\qquad
\forall u=(\operatorname{curl}\psi,\omega),
\quad v=(\operatorname{curl}\phi,\theta)\in V,
\]
$V$ 是 Hilbert 空间, 且 Stokes 算子的定义可延拓为
\[
T\in\mathcal L(X';V),
\qquad
\widetilde a(Tl,v)=\langle l,v\rangle
\quad\forall v\in V.
\]
由 \eqref{eq:navier-centered-stream-adjoint-derivative} 立即得到
\[
\widetilde a(TD^*z,v)=\widetilde a(TDv,z)
\qquad\forall z,v\in V.
\]
换言之, 对标量积 $\widetilde a(\cdot,\cdot)$,
$TD^*$ 是 $TD$ 在 $V$ 中的伴随. 因而, 由于假设 $I+TD$ 是 $V$ 上的同构,
$I+TD^*$ 也同样是 $V$ 上的同构.
